(1) 设 \(x(t)=A(t)e^{i\omega_0t}+A^*(t)e^{-i\omega_0t}\)，慢幅度满足 \(|\dot A|\ll\omega_0|A|\)，忽略 \(\ddot A\) 等二阶小量。参数调制中的 \(\cos(2\omega_0t)\) 与共轭快因子相乘会产生共振的 \(e^{i\omega_0t}\) 分量。

(2) 只保留与 \(e^{i\omega_0t}\) 共振的项，得到
\[
2i\omega_0\dot A+2i\gamma\omega_0A+\frac{h\omega_0^2}{2}A^*=0,
\quad
\dot A=-\gamma A+i\frac{h\omega_0}{4}A^*.
\]
复共轭式为 \(\dot A^*=-\gamma A^*-i(h\omega_0/4)A\)，是二维实线性慢变量系统。

(3) 该矩阵的两个本征增长率为
\[
\lambda_\pm=-\gamma\pm\frac{|h|\omega_0}{4}.
\]
于是当 \(|h|>4\gamma/\omega_0\) 时 \(\lambda_+>0\)，主参数不稳定；反之两个慢模都衰减。恰在阈值时领先一阶近似给零增长率，细致判断需更高阶项。

(4) 若 \(\gamma=0\) 而 h 非零，任意足够小的精确共振调制均在领先阶有指数增长；若 h=0，则自由阻尼振子按 \(e^{-\gamma t}\) 衰减。指数增长最终使 x 不再“小”，实际装置中的几何非线性、阻尼变化或外部能量饱和会终止该近似，不能从线性方程断言物理振幅无限大。
